\documentclass[a4paper]{article}

\usepackage{graphicx}
\usepackage{amsmath,amssymb}
\usepackage{minted}
\usepackage{multirow}
\usepackage{float}
\usepackage{todonotes}
\usepackage{forest}
\usepackage{xstring}
\usepackage{xcolor}
\usepackage[most]{tcolorbox}
\usepackage{xparse}
\usepackage{natbib}

\usetikzlibrary{graphs,graphdrawing}
\usegdlibrary{layered}

% for external graphviz
\input{/home/donkarlo/Dropbox/repo/nd_language_ai_project/src/nd_language_ai/formal/computer/programming/tex/latex/code_clips/external_graphviz.tex}

% for tags
\input{/home/donkarlo/Dropbox/repo/nd_language_ai_project/src/nd_language_ai/formal/computer/programming/tex/latex/parts/control_sequence/kind/semtag/semtag.tex}

% for links
\input{/home/donkarlo/Dropbox/repo/nd_language_ai_project/src/nd_language_ai/formal/computer/programming/tex/latex/code_clips/seehere.tex}

\title{}
\date{}
\author{Mohammad Rahmani}

\begin{document}

    \maketitle

    \section{Thema: Sollten Menschen weniger soziale Medien benutzen?}

        Schreibe etwa 80--100 Wörter. Du kannst diese Punkte benutzen:

        \begin{itemize}
            \item Wie oft benutzt du soziale Medien?
            \item Welche Vorteile haben sie?
            \item Welche Nachteile gibt es?
            \item Was ist deine persönliche Meinung?
            \item Was würdest du anderen Menschen empfehlen?
        \end{itemize}
        Mit `%` geht das **nicht innerhalb derselben Zeile**, weil LaTeX nach `%` den gesamten Rest der Zeile ignoriert. Wenn du **keine neue Zeile für jede Klammer** willst, können wir stattdessen `\iffalse ... \fi` verwenden. Der Inhalt wird von LaTeX vollständig ignoriert, bleibt aber genau an derselben Stelle im Text.

        ```latex
        Persönlich benutze ich viel soziale Medien.

        Ich denke, dass die sozialen Medien sehr viel sinnvoll für die Enwicklung der Demokratie sein können.

        Ich weiß, dass wir bisher\iffalse so far\fi es nicht bezeugt\iffalse to withness\fi haben, aber wir müssen auch beachten, dass\iffalse dass spielst der Objektroll weil beaten brauchst ein davon(yeki az an=yek objekt)\fi, obwohl diese Medien nicht sehr jung sind, sie auch nicht genug erwachsene\iffalse Adult\fi sind, damit wir erwarten, dass sie über unsere democratikche Zukunft ein effektiv Effekt zu haben. Aber, trotz allem, was ich gessagt habe, bezeugen\iffalse to witness\fi wir, dass diese Medien den stimmlosen Völker ein Mittle, um sich auszudrüken.

        Da ich denke, dass jedes Massenmedium\iffalse singular neutral\fi, das\iffalse pronom für Massenmedien\fi der Gesellschaft\iffalse f\fi hinzufügt wird \iffalse von Behörde\fi, dem Volk hift, durch Narrichten besser infrormiert zu werden.

        Nun hat sich die Frage in die verwandelt\iffalse to change, benutzt immer mit "in". Es heißt wir benutzen immer verwandeln in etwas die hier wir "in" dahin benutzt\fi, wie die Narrichten uns helfen, auf eine bessere Demokratie\iffalse f\fi zuzumarschieren \iffalse Die zweite zu ist die zu die ist für infinitiv bauen benutzt worden\fi.

        Nachrichten helfen uns, uns\iffalse "uns" ist eine Variante von Refelxivpronomen für sich informieren die später in deselb Satz benutzt wird\fi besser über die Strecke\iffalse m, dativ Objekt\fi die, wir nehmen müssen, zu informieren \iffalse Das verb ist "sich informieren" aber "uns" ist vor \fi, um unsere Rechte\iffalse plural hier aber Recht is neutral\fi besser zu verteidigen.

        Sogar vor\iffalse before, andres als bevor die Hauptsatz-nebensatz Konstruktion nötigt, vor muss immer vor ein Nom kommen. Weil es kein Bewegung gibt, dann Mann muss Dativ benuzen\fi der Strecke, die wir nehmen müssen, um unsere Rechte zu verteidigen, müssen wir über sie uns informieren.

        Außerdem müssen wir entdecken\iffalse to discover\fi, welche politische Partei uns hilft \iffalse hilfen verb steht sich am ende, weil es "Welche" hat ein neben satz gebaut\fi, \iffalse kein um...zu benutz wird, weil das "Helfen" konstruktion folgt "jemand Hefen zu etwas machen". Wir müssen wissen, obwohl was in die letze Phrase erklärt worden ist man kann helfen mit um ... zu benutzen aber die Bedutung wird ein bischen verschiden\fi unsere Rechte besser zu verteideigen.

        Wie es bisher erklärt worden ist\iffalse Nebenstaz\fi, brauchen wir \iffalse zweichendrin die \fi die richtigen Nachrichten\iffalse feminin plural\fi, um eine praktische DemoKratie\iffalse f\fi zu haben.

        Wir können entweder ein einziges\iffalse wegen ein und Medium, wir brauchen ein gemischte Adjektiv, deshalb wir -es hinzufügen\fi Medium,\iffalse Neutrum Akkusativ Singular\fi von dem wir richtigesten Nachrichten\iffalse Akkusativobjekt von erwarten\fi erwarten, oder mehrere Medien von denen \iffalse Dativ plural\fi, das Publikum die Nachrichten vergleichen kann, um die präzisesten Nachrichten\iffalse Akkusativ plural für entnehmen\fi zu entnehmen\iffalse to extract\fi.

        Ich verteidige mehrere Medien, weil ich finde überhaupt nicht möglisch, dass wir von ein einzige sozial Medium erwarten, dass es ganz neutral bleibt.

        Deshlab\iffalse Deshalb macht nicht der Satz ein Nebensatz\fi finde ich mehrere Medien als die Quellen, von denen\iffalse Dativ weil es nachg von stelt und feminin, plural weil es ist ein Pronom für Quelen\fi wir über unsere Rechte\iffalse die Plural ist Rechte und der Nom Konjugirt nicht mit andren Teile des Satz \fi und die politische Partei, die unsere Rechte verteidigt lernen können.

        Soziale Medien\iffalse pl\fi verbreiten\iffalse AKK obj\fi uns die außergewöhnlichen\iffalse Schwach\fi Möglichkeiten\iffalse Plural Nom\fi um jeden\iffalse es ein Pronom hier ist und laut umwandeln, es umwandelt sich in Akkusativ Objekt. Es ist kein Artikel\fi individuell nach einen Jornalist\iffalse m\fi umzuwandeln\iffalse to convert. umwandeln etwas in etwas\fi.

        Nach den\iffalse Plural Dativ, weil es ein Dativ objekt für verbinden ist\fi obigen\iffalse Schwach Adjektiv\fi Argumenten\iffalse PluralDativ, weil normalweise nach ``nach`` kommt ein Dativ\fi\iffalse hier braichen wir kein komma, weil Nach den obigen ... ist der Objekt die an der Anfang gekommt hat statt der Subjekt, die Soziale Medien ist\fi sind die sozialen\iffalse Schwach nominativ Adjektiv für Medien, die der Subjekt des Satzs ist\fi Medien dierkt mit \iffalse Dativ bestimmte Artikel\fi Extraktion\iffalse f\fi der\iffalse Genitiv\fi richtigen\iffalse Adjetiv Genetiv\fi Nachrichten\iffalse Plural\fi verbunden, auf denen\iffalse Relativpronom\fi Demokratie \iffalse f, Dativ\fi basiert.
        ```

        So bleiben **Absätze und Zeilenstruktur erhalten**, und jede frühere Klammer steht direkt dort, wo sie vorher stand.


        \bibliographystyle{plainnat}
        \bibliography{/home/donkarlo/Dropbox/repo/nd_shared_research/bibliography/bibliography.bib}

\end{document}